\section{Related Work}

\noindent\textbf{Patent-Paper Pairs (PPPs). }
Many research labs practice concurrent patenting and academic publishing.
\citet{murrayFormalIntellectualProperty2005} find that almost 50\% of their sampled academic papers from \textit{Nature Biotechnology} have a corresponding US patent.
In economics, PPPs have been used to study innovation dynamics, like whether patenting promotes or hinders the free flow of innovations \cite{murrayFormalIntellectualProperty2005,magermanSearchAnticommonsPatentpaper2011}.
In that context, several approaches to finding PPPs have been proposed:
\citet{murrayInnovationCoevolutionScientific2002} and \citet{murrayFormalIntellectualProperty2005} identify pairs manually by analyzing their full texts and citation networks.
\citet{magermanExploringFeasibilityAccuracy2010} and \citet{vanlooyAssessmentLatentSemantic2011} explore several data mining techniques.
We refine and extend their approach, and add criteria for author overlaps, date ranges, competing candidates, and licenses.
\citet{gansPatentsPapersPairs2008} propose a taxonomy of PPPs that includes both to 1-to-1 matches and m-to-n matches.
To ensure that papers are a solid source of information about the invention, we design our matching procedure to find only 1-to-1 matches.
We publish, to the best of our knowledge, the first PPPs dataset for NLP research.


\begin{figure*}
  \small
  \setlength{\tabcolsep}{10pt}
  \begin{tabularx}{\textwidth}{p{.95\textwidth} p{0\textwidth}}
    \centering
    \begin{tabular}[t]{ccccccc}
      \toprule
      \multirow{2}{*}{\textbf{Split}} & \multirow{2}{*}{\textbf{\# pairs}} & \multirow{2}{*}{\textbf{\# patent tokens}} & \multirow{2}{*}{\textbf{\# paper tokens}} & \multicolumn{3}{c}{\textbf{\# outline bullets}}                                       \\
      \cmidrule(lr){5-7}
                                      &                                    &                                            &                                           & short                                           & medium          & long              \\
      \midrule
      train                           & 1000                               & 17.8k $\pm$ 15.1k                      & 7.9k $\pm$ 4.5k                       & 36.8 $\pm$ 29.2                                 & 73.5 $\pm$ 60.0 & 149.0 $\pm$ 122.0 \\
      val                             & 242                                & 18.2k $\pm$ 16.3k                      & 8.0k $\pm$ 4.1k                       & 37.4 $\pm$ 30.8                                 & 74.4 $\pm$ 62.9 & 150.6 $\pm$ 127.5 \\
      test                            & 500                                & 18.1k $\pm$ 13.2k                      & 8.1k $\pm$ 3.9k                       & 37.5 $\pm$ 24.7                                 & 74.9 $\pm$ 50.9 & 151.6 $\pm$ 103.7 \\
      nc-test                         & 71                                 & 18.4k $\pm$ 13.9k                      & 9.5k $\pm$ 4.6k                       & 37.8 $\pm$ 22.7                                 & 76.0 $\pm$ 46.4 & 154.4 $\pm$ \phantom{0}94.8  \\
      \midrule
      all                             & 1813                               & 17.9k $\pm$ 14.7k                      & 8.1k $\pm$ 4.3k                       & 37.1 $\pm$ 28.0                                 & 74.1 $\pm$ 57.5 & 150.1 $\pm$ 117.0 \\
      \bottomrule
    \end{tabular}
    \captionof{table}{\DatasetName statistics. Values are reported as mean $\pm$ std. Token counts correspond to the Llama-3 tokenizer.}
    \label{tab:dataset_stats} &
  \end{tabularx}
\end{figure*}



\noindent\textbf{Patent Generation. }
Most prior work on patent generation has focused on titles, abstracts, and claims.
\citet{christofidellisPGTPromptBased2022} train a multi-task GPT2 model to generate these parts.
\citet{leeEvaluatingGenerativePatent2023} pre-trains a GPT-J-6B architecture on entire patents and evaluates it on claim generation.
\citet{zuoPatentEvalUnderstandingErrors2024} use GPT-3.5-turbo and Llama-2 to generate abstracts from claims, and claims from previous claims.
\citet{wangPatentformerNovelMethod2024} experiment with the generation of individual description paragraphs based on patent claims and drawing descriptions.
In practice, it is not likely that claims are already finalized at the time of writing of the description section.
In their work concurrent to ours, \citet{wangAutoPatentMultiAgentFramework2024} leverage an agent framework to generate complete patents based on purely automatically generated invention specifications.
They share our core principle of using a divide-and-conquer strategy for long-document generation, but evaluate only using surface-level metrics.
Our more realistic setting relies on real-world invention specifications and provides deeper insights due to more sophisticated evaluation metrics.

\noindent\textbf{Outline-guided Generation}
is a paradigm in which LLMs use an outline to produce longer, more structured and coherent text.
Outlines are either generated in a planning stage 
\cite{drissiHierarchicalTextGeneration2018,yaoPlanandWriteBetterAutomatic2019,sunSummarizeOutlineElaborate2022a,yang-etal-2022-re3,leeNavigatingPathWriting2024,wangAutoSurveyLargeLanguage,shaoAssistingWritingWikipedialike2024} 
or provided as input 
\cite{fangOutlineStoryFinegrained2021a,spangher-etal-2022-sequentially,yangDOCImprovingLong2023,liAdvancingPreciseOutlineConditioned2023}.
They are created using extraction of key words 
\cite{yaoPlanandWriteBetterAutomatic2019}, 
phrases 
\cite{fangOutlineStoryFinegrained2021a}, 
or sentences 
\cite{drissiHierarchicalTextGeneration2018,sunSummarizeOutlineElaborate2022a,yang-etal-2022-re3,liAdvancingPreciseOutlineConditioned2023}, 
generated using LLMs 
\cite{yangDOCImprovingLong2023} or 
defined interactively 
\cite{goldfarb-tarrant-etal-2019-plan}.
In our work, we posit that outlines should be provided by the patent attorney to satisfy the high demand for user control. %

